\chapter{Continual Learning (Anti-forgetting)}

\ChapterMeta{
This chapter covers continual fine-tuning across tasks and domains, and the evaluation of forgetting using recorded run artifacts.
}{
It helps mitigate catastrophic forgetting and makes per-task metrics comparable over a training sequence.
}{
It requires continual configs and run artifacts (for example \cmd{continual_run.json}); optional replay and LoRA settings change compute and forgetting behavior.
}
\ChapterAudience{Read this chapter when you are training across tasks or domains and need to reason about forgetting, replay, self-distillation, LoRA, or QLoRA.}


\section{What YOLOZU means by ``continual''}
Continual learning (CL) here means fine-tuning across a sequence of tasks/domains while mitigating catastrophic forgetting.
In this repository, the reference implementation targets the in-repo \path{rtdetr_pose} scaffold.

\section{Core approach (baseline)}
The default strategy combines:
\begin{itemize}
  \item \textbf{Memoryless self-distillation}: keep the new model close to the previous checkpoint while learning the new task.
  \item \textbf{Optional replay buffer}: add a small buffer of past samples (e.g., reservoir sampling) and train on \emph{(current task + replay)}.
  \item \textbf{Optional parameter-efficiency}: restrict trainable parameters (e.g., LoRA) to reduce forgetting pressure.
\end{itemize}

\begin{figure}[h]
\centering
\resizebox{\linewidth}{!}{%
\begin{tikzpicture}[
  font=\small,
  box/.style={draw, rounded corners, align=left, inner sep=6pt},
  arrow/.style={-{Latex[length=2mm]}, thick}
]

\node[box] (tasks) at (0,0) {\textbf{Incoming task stream}\\Task A $\rightarrow$ Task B $\rightarrow$ Task C};
\node[box] (prev) at (4.8,0.9) {\textbf{Previous checkpoint}\\acts as a teacher / anchor};
\node[box] (replay) at (4.8,-0.9) {\textbf{Replay buffer}\\small memory of older samples};
\node[box] (train) at (9.5,0) {\textbf{New training run}\\current task + replay + self-distill};
\node[box] (eval) at (14.0,0) {\textbf{Evaluation}\\old tasks and new task are measured together};

\draw[arrow] (tasks) -- (train);
\draw[arrow] (prev) -- (train);
\draw[arrow] (replay) -- (train);
\draw[arrow] (train) -- (eval);

\end{tikzpicture}
}
\caption{Continual learning in \yolozu{}: the new run does not train on the latest task alone; it is anchored by the previous checkpoint and may also replay a small memory of older examples.}
\end{figure}

\section{LoRA / QLoRA in plain language}
LoRA and QLoRA are not separate continual-learning algorithms.
They are \textbf{ways to make the parameter updates smaller and cheaper}.

\subsection{LoRA}
The easiest mental model is:
\begin{quote}
freeze most of the big model, then train a much smaller correction on top.
\end{quote}

That is why LoRA is attractive in continual learning:
\begin{itemize}
  \item a smaller trainable surface can reduce forgetting pressure,
  \item it is often easier to fit within a limited GPU budget,
  \item it gives a more controlled ablation than full-model fine-tuning.
\end{itemize}

\subsection{QLoRA}
QLoRA keeps the same adapter idea, but stores the frozen base model in a more memory-efficient quantized form.

In plain language:
\begin{itemize}
  \item \textbf{LoRA}: freeze the large model, train small adapters.
  \item \textbf{QLoRA}: do the same thing, but keep the frozen base cheaper to hold in memory.
\end{itemize}

\begin{figure}[h]
\centering
\resizebox{\linewidth}{!}{%
\begin{tikzpicture}[
  font=\small,
  box/.style={draw, rounded corners, align=left, inner sep=6pt, text width=0.27\linewidth},
  arrow/.style={-{Latex[length=2mm]}, thick}
]

\node[box] (base) at (0,0) {\textbf{Full fine-tuning}\\All main weights update.\\Highest flexibility, highest memory / forgetting pressure.};
\node[box] (lora) at (6.3,0) {\textbf{LoRA}\\Base model stays mostly frozen.\\Small low-rank adapters carry the update.};
\node[box] (qlora) at (12.6,0) {\textbf{QLoRA}\\Same adapter idea as LoRA, but the frozen base is kept in a quantized form for lower memory use.};

\draw[arrow] (base.south) -- ++(0,-0.9) node[below, align=center] {largest update surface};
\draw[arrow] (lora.south) -- ++(0,-0.9) node[below, align=center] {smaller, easier-to-control update};
\draw[arrow] (qlora.south) -- ++(0,-0.9) node[below, align=center] {adapter update + cheaper frozen base};

\end{tikzpicture}
}
\caption{LoRA and QLoRA are best understood as parameter-efficiency tools. They change \emph{how much} of the model moves during continual learning, not the underlying anti-forgetting objective itself.}
\end{figure}

\subsection{Pros and Cons}
\textbf{LoRA}
\begin{itemize}
  \item Pros: simpler mental model than QLoRA; often cheaper than full fine-tuning; integrates naturally with replay and self-distillation.
  \item Cons: still introduces extra knobs; can underfit if the adapter is too small; may hide the fact that the base model is the real bottleneck.
\end{itemize}

\textbf{QLoRA}
\begin{itemize}
  \item Pros: lowers memory pressure further; keeps the adapter-style workflow; attractive when the frozen base dominates footprint.
  \item Cons: adds quantization complexity; depends on backend support; is more experimental in this repository than plain LoRA.
\end{itemize}

\subsection{A practical choice example}
Imagine a domain-incremental warehouse-camera project:
\begin{itemize}
  \item the previous checkpoint is already usable,
  \item the main concern is forgetting older camera views,
  \item GPU memory is limited, but not critically constrained.
\end{itemize}

In that situation, \textbf{LoRA is usually the better first experiment}.
It keeps the story simpler: replay + self-distillation still do the anti-forgetting work, while the adapter merely reduces how much of the model moves.

Now change one assumption:
\begin{itemize}
  \item the frozen base is the real memory bottleneck,
  \item batch size collapses unless the base becomes cheaper to hold.
\end{itemize}

That is when \textbf{QLoRA becomes the practical option}.
You keep the same adapter-style update logic, but trade some simplicity for a smaller memory footprint.

So the operational rule is:
\begin{itemize}
  \item \textbf{LoRA first when clarity and debugging speed matter},
  \item \textbf{QLoRA when memory is the real blocker}.
\end{itemize}

\section{Main entry points}
\subsection{Train}
Run continual fine-tuning using a config (start from the example):
\begin{lstlisting}[language=bash]
python3 rtdetr_pose/tools/train_continual.py \
  --config configs/continual/rtdetr_pose_domain_inc_example.yaml
\end{lstlisting}

To run \textbf{memoryless} (no replay), set replay size to 0 in the config.

\paragraph{Note on "distillation" terminology}
This chapter's \emph{distillation} refers to \textbf{checkpoint-based self-distillation during training} (e.g., \cmd{--self-distill-from}) as an anti-forgetting mechanism.
The separate tool \path{tools/distill_predictions.py} performs \textbf{offline prediction blending} (rewriting a predictions JSON) and does not by itself mitigate catastrophic forgetting.

\subsection{Evaluate}
Evaluate forgetting and per-task metrics from the recorded run JSON:
\begin{lstlisting}[language=bash]
python3 tools/eval_continual.py \
  --run-json runs/continual/<run>/continual_run.json \
  --device cpu \
  --max-images 50
\end{lstlisting}

If your dataset provides pose sidecar metadata (e.g., under \path{labels/<split>/*.json}), you can request pose metrics:
\begin{lstlisting}[language=bash]
python3 tools/eval_continual.py \
  --run-json runs/continual/<run>/continual_run.json \
  --device cpu \
  --max-images 50 \
  --metric pose \
  --metric-key pose_success
\end{lstlisting}

Available \cmd{--metric-key} values in pose mode:
\begin{itemize}
  \item \cmd{pose_success}, \cmd{rot_success}, \cmd{trans_success}, \cmd{match_rate}, \cmd{iou_mean}
  \item \cmd{depth_abs_mean}, \cmd{depth_abs_median}
  \item \cmd{add_mean}, \cmd{add_median}, \cmd{adds_mean}, \cmd{adds_median}
\end{itemize}

Notes:
\begin{itemize}
  \item Depth metrics require GT depth/translation and predicted depth/translation fields.
  \item ADD/ADDS metrics require CAD points in sidecar metadata (\cmd{cad_points} / \cmd{cad_path} / \cmd{cad_points_path}) and valid GT+pred pose.
\end{itemize}

Examples for metric-key switching:
\begin{lstlisting}[language=bash]
# depth error summary
python3 tools/eval_continual.py \
  --run-json runs/continual/<run>/continual_run.json \
  --metric pose \
  --metric-key depth_abs_mean

# ADD / ADDS (requires CAD points)
python3 tools/eval_continual.py \
  --run-json runs/continual/<run>/continual_run.json \
  --metric pose \
  --metric-key add_mean

python3 tools/eval_continual.py \
  --run-json runs/continual/<run>/continual_run.json \
  --metric pose \
  --metric-key adds_mean
\end{lstlisting}

\section{Run artifacts (what to version / what to compare)}
\cmd{train_continual.py} writes a run folder under \path{runs/continual/}.
Key outputs:
\begin{itemize}
  \item \path{continual_run.json}: task list, checkpoints, config hash, and run record.
  \item \path{replay_buffer.json}: buffer summary (when replay is enabled).
  \item Per-task subfolders, typically including:\\
    \path{checkpoint.pt}, \path{metrics.jsonl}, \path{metrics.csv}, \path{run_record.json}.
\end{itemize}

For reporting and plots, treat \path{continual_run.json} as the single source of truth.

\section{Selected config knobs (high leverage)}
Commonly tuned items include:
\begin{itemize}
  \item Replay: \cmd{replay_size}, \cmd{replay_fraction}, \cmd{replay_per_task_cap}.
  \item Distillation: \cmd{distill.enabled}, \cmd{distill.keys}, \cmd{distill.weight}, \cmd{distill.temperature}.
  \item Replay distillation (DER++-style): \cmd{derpp.enabled} and related keys (requires teacher outputs stored in records).
  \item Regularizers across tasks (optional): EWC (\cmd{ewc.*}) and SI (\cmd{si.*}).
\end{itemize}

\section{Research workflow guidance}
\begin{itemize}
  \item \textbf{Pin the evaluation subset}: use a deterministic subset (same images each run) before comparing CL runs.
  \item \textbf{Log runtime cost}: replay and distillation change throughput; report accuracy vs cost.
  \item \textbf{Keep baselines simple}: compare (A) naive fine-tune, (B) memoryless self-distill, (C) replay+distill.
\end{itemize}

\section{Notes}
\begin{itemize}
  \item Some continual comparisons use a CPU-friendly proxy metric (e.g., a lightweight mAP-style summary) on a pinned subset to keep iteration fast.
  \item For ``real'' mAP, switch your evaluation workflow to a full COCO evaluator when appropriate.
\end{itemize}

\section{References (self-distillation background)}
This repository's checkpoint-based self-distillation is a pragmatic anti-forgetting regularizer.
For background reading:
\begin{itemize}
  \item Knowledge distillation (original): \cite{hinton2015distilling}.
  \item Self-distillation (classic reference): \cite{furlanello2018bornagain}.
  \item Continual learning via self-distillation (SDFT): \cite{shenfeld2026sdft}.
  \item LoRA as parameter-efficient tuning: \cite{hu2021lora}.
  \item QLoRA as quantized adapter tuning: \cite{dettmers2023qlora}.
\end{itemize}
